% RESEARCH_DRAFT_SHA256: 2ca893947d93337f8aa111d3702c0ca7f11671d5231070a0867aefaa34cd0d47
% Layout and source-reference wrappers expanded for the web renderer.
\documentclass[11pt]{article}
\usepackage[margin=25mm]{geometry}
\usepackage{fontspec,unicode-math}
\setmainfont{Latin Modern Roman}
\setmathfont{Latin Modern Math}
\usepackage{amsmath,amssymb,mathtools,xurl,hyperref}
\hypersetup{colorlinks=true,urlcolor=blue,linkcolor=blue}
\setcounter{secnumdepth}{0}
\setlength{\parindent}{0pt}
\setlength{\parskip}{5pt}
\setlength{\emergencystretch}{3em}
\newcommand{\sref}[2]{\hyperlink{source:#1:#2}{[S#2]}}
\newcommand{\cataloguescope}[1]{\paragraph{Source record.}#1}
\newcommand{\floor}[1]{\left\lfloor#1\right\rfloor}
\begin{document}
\section*{Research-attempt provenance}
This standalone research attempt retains the complete original definition body below,
including its historical source and preparation notes. Those notes describe the
earlier catalogue preparation; the new attempt and its verification are identified
separately after that body. No original ID, catalogue status, or source statement is changed.

% UnsolvedMath ID 20001206; source catalogue code AIM-COMPUTATION-0044
\section{Vertex flattenings generate the phylogenetic ideal}\label{20001206}
{\small UnsolvedMath ID 20001206 · Computer Science}
\subsection{Definitions and mathematical statement}

\paragraph{Shared notation.}Standard number systems are $\mathbb N,\mathbb Z,\mathbb Q,\mathbb R,\mathbb C$. Specialized notation and index ranges follow the cited source and the definitions below.


A phylogenetic model assigns a probability distribution to observed leaf states on a tree using hidden internal states and stochastic transition matrices along edges. A phylogenetic invariant is a polynomial vanishing on that model. Flattening groups tensor indices into blocks to produce a matrix or lower-order tensor; rank constraints can then give polynomial minors. The exact constraints depend on the mutation model and state cardinality.
\paragraph{Record scope and qualifications.}
The catalogue title is retained as metadata. Named certificates, contracts, bounds or frameworks in that title are not asserted as results of the original source extract.
The tree model, state size, and meaning of three-dimensional flattening are not supplied in the excerpt. These must be identified before stating a precise ideal equality; the catalogue marks the historical problem solved. The supplied catalogue marks this record solved; its present research status was not independently checked.
The source text below is rendered literally where necessary; damaged formulas and composite extracts are not silently repaired into a new mathematical claim.
\paragraph{Statement recorded in the supplied source.}
Question 30. Is the ideal obtained by taking the invariants arising from all 3-dimensional flattenings equal to the ideal of all invariants for a phylogenetic tree? \href{https://aimath.org/WWN/compalgstat/compalgstat.pdf}{[S2]}
\paragraph{Source status.}
The supplied catalogue labels this source record solved; that metadata was checked against the supplied snapshot on 2026-10-11. The label is preserved for provenance and is not an independent verification of present open status or mathematical correctness. This is documentation of a source-labelled solved record, not a newly proposed unsolved problem.
\subsection{Short English statement}
The historical question asks whether invariants from all three-way flattenings generate the full phylogenetic invariant ideal.
\subsection{Sources}
\begin{itemize}
\item[S1] Ulam AI and the UnsolvedMath Contributors, UnsolvedMath: A Curated Collection of Open Mathematics Problems, record 20001206 (AIM-COMPUTATION-0044). Catalogue curation and metadata: CC BY 4.0. The source extract is retained; definitions, English summaries and qualifications are editorial additions.
\url{https://huggingface.co/datasets/ulamai/UnsolvedMath}.
\item[S2] AIM Workshop Problem Lists, source linked by catalogue record AIM-COMPUTATION-0044.
\url{https://aimath.org/WWN/compalgstat/compalgstat.pdf}.
\end{itemize}
\paragraph{Preparation and review.}
The original source and UnsolvedMath contributors retain attribution for the source extract. Codex (GPT-6; exact serving revision not exposed) assisted with definitions, English summaries, scope notes and formatting. No new human mathematical authorship or independently verified solution is asserted. Permanent source ID 20001206 and code AIM-COMPUTATION-0044 checked against the complete supplied registry and current repository. Distinct upstream records, including repeated questions and broad or damaged extracts, are retained for complete coverage of the source CS category. A repeated record is not a new independent conjecture. Complete mathematical-equivalence and current-literature review has not been performed. Preparation and final review follow contribution prompts C and D, with source gaps disclosed. Structural and rendering checks do not establish a new theorem or open status.
\label{end:20001206}

\section{Literature and context check}
The original AIM question and Draisma--Kuttler's primary arXiv abstract,
\emph{On the ideals of equivariant tree models}, were consulted. The latter
reports generation of the full tree ideal by vertex-flattening ideals in its
framework. Its full proof and all model hypotheses were not independently
checked here. This attempt preserves the catalogue's solved status and examines
two invalid shortcuts; it does not label the historical theorem newly open.
\section{Research attempt: Matrix rank tests cannot replace full vertex-flattening ideals}
A proposed shortcut interprets every three-way flattening merely through
ordinary matrix ranks. Test this on a $3\times3\times3$ tensor with slices
$A,B,C$. For a rank-at-most-three decomposition
$T=\sum_{j=1}^3u_j\otimes v_j\otimes w_j$, write
$A=U\operatorname{diag}(w_{1j})V^T$, and similarly for $B,C$.
If $A$ is invertible, then $U,V$ and its diagonal factor are invertible and
direct multiplication gives
\[
 BA^{-1}C=CA^{-1}B,
 \qquad B\operatorname{adj}(A)C-C\operatorname{adj}(A)B=0.
\]
The latter is a polynomial identity in the tensor entries. It holds also for
all rank-at-most-three parameter choices with singular $A$: approximate their
$U,V,w$ by choices with invertible $U,V$ and all $w_{1j}\ne0$, use the identity
there, and take the limit. Invertible matrices are dense because
$\det(U+tI)$ is a nonzero polynomial in $t$, and similarly for $V$.
It therefore also holds on the closure of the rank-three parameterized set.

Choose $A=I_3$, $B=E_{12}$, $C=E_{21}$. Then the polynomial matrix equals
$\operatorname{diag}(1,-1,0)$ and does not vanish. Yet every ordinary
flattening is a $3\times9$ matrix or its transpose, automatically of rank at
most three. Thus those matrix conditions cannot characterize even this small
three-way hidden-class model. Full vertex-flattening ideals contain information
not provided by the tautological matrix-rank tests.

Another proposed shortcut proves only equality of zero sets. In
$\mathbb C[x,y]$, the ideals $I=(x)$ and $J=(x^2,xy)$ have the same zeros,
the line $x=0$, but $x\notin J$: every nonzero homogeneous component of an
element of $J$ has degree at least two. Hence $J\ne I$. This exact example
shows that set-theoretic reconstruction alone would not answer the source's
ideal-equality question.

Neither obstruction refutes generation by the full ideals at vertices. The
catalogue already labels the historical question solved, and the consulted
Draisma--Kuttler abstract reports precisely full ideal generation in its
equivariant-tree framework. That literature is context, not a theorem newly
proved by the two elementary diagnostics above.
\section{Verification and adversarial audit}
The noncommuting slices and commutator diagonal were checked by multiplying
matrix units. Singular slice parameters are covered by an explicit density
argument, not by dividing by a zero determinant. The ideal example checks
equality of zero sets separately from failure of ideal membership.
\paragraph{Audit corrections.} The source solved label is preserved, full ideal generation is not confused with set equality, and matrix flattenings are not silently substituted for three-way vertex ideals.
\paragraph{External computation.} None; symbolic arguments and boundary checks only.
\section{Scope and remaining gaps}
No new result on the historical solved theorem is claimed. The source
extract omits its precise model and field; a full proof of the appropriate
literature theorem is not reproduced. The investigated shortcut fails because
it substitutes matrix ranks or radicals for full vertex-flattening ideals.
\paragraph{Integration review.} Codex (GPT-6; exact serving revision not exposed) reviewed the written arguments and expanded layout and source-reference wrappers. Repository metadata and the production web renderer were checked using Python 3.13 and JavaScript with MathJax 3.2.2. This model-assisted review is not independent human proof verification; no PDF compilation or formal proof checking is asserted.

\section{Completion Estimate}
\noindent\textbf{COMPLETION: 0\%}
The calculations below the preserved source statement investigate a concrete route,
but establish no substantiated new progress toward the entire original target.
Known special cases and failed sufficient conditions are not credited as a solution.
\section{References}
The original source references and catalogue attribution are preserved in the
Sources section above. The following URLs were actually consulted for this attempt;
retention of an additional original reference does not claim fresh consultation of it.
\begin{itemize}
\item \url{https://aimath.org/WWN/compalgstat/compalgstat.pdf} (consulted 2026-10-11). 
\item \url{https://arxiv.org/abs/0712.3230} (consulted 2026-10-11). 
\end{itemize}
\section{Model and contributor disclosure}
OpenAI gpt-6-astra, requested reasoning configuration ultra, generated this research
attempt on 2026-10-11. The exact serving revision is not exposed. Codex GPT-6
performed orchestration, repository-format review, and a model-assisted mathematical
audit of the written arguments; its exact serving revision
is also not exposed. The mathematical route, boundary checks and final adversarial
self-audit followed contribution prompts A then B in
\texttt{docs/contribution\_prompts.md}, after reading \texttt{CONTRIBUTING.md}.
Python 3.14.2 assembled the standalone documents; any mathematical computation is
specified in Verification. No theorem prover was used. No human mathematical
authorship, substantial human proof correction, or independent human proof checking
is asserted. Model-assisted review, source retrieval, and structural checks do not
establish correctness or novelty. Full private reasoning transcripts are unavailable
as reproducibility artifacts; the complete arguments, inputs and outcomes actually
used are stated here or in the companion batch audit.
\section{Final status}
Unresolved as a new research contribution. This is a concrete route investigation,
not a claim of new complete resolution or a current-status certification. A known
answer reproduced for context does not become a new contribution. The catalogue's
source-status label is preserved; this attempt status does not assert that a
source-labelled solved question is open.
\end{document}
